\usepackage{tikz}
\usepackage{tikzscale}
% adapted from https://stackoverflow.com/a/54903109/1305099
\NewDocumentCommand{\semitransp}{m}{%
  {\pgfsetfillopacity{0.3}#1}%
  \pgfsetfillopacity{1}%
}
\usepackage{graphicx} % support the \includegraphics command and options

\usepackage{multirow}%for tables

% Font/encoding setup, guarded by engine: under pdfLaTeX [T1]{fontenc} is
% required or monospace silently falls back to a proportional font.  didactic's
% memoir branch overrides the main fonts (Palatino/Helvetica/Bera); lmodern is
% the fallback for the beamer build.
\usepackage{iftex}
\ifpdftex
  \usepackage[utf8]{inputenc}
  \usepackage[T1]{fontenc}
  \usepackage{lmodern}
\fi
\usepackage[british]{babel}
\usepackage{booktabs}

% Verbose citation style (full reference on first use, short form on repeats):
% \autocite produces a footnote citation, which didactic places in the margin in
% the memoir article and at the foot of the slide in beamer -- "verbose
% references in the margin".  Use \autocite for all citations so they adapt to
% the build; keep \parencite/\textcite out of the document.
\usepackage[
  natbib,
  backend=biber,
  style=verbose,
  citestyle=verbose,
  singletitle=false,
  maxbibnames=99,
]{biblatex}
\addbibresource{debugging.bib}
\addbibresource{theory.bib}

\usepackage[all]{foreign}
\renewcommand{\foreignfullfont}{}
\renewcommand{\foreignabbrfont}{}

%\usepackage{newclude}
\usepackage{import}

\usepackage{subcaption}

\usepackage[noend]{algpseudocode}
\usepackage{xparse}

\let\email\texttt

% minted v3+ (TeX Live 2024+) auto-detects the output directory; the old
% [outputdir=ltxobj] option is now an error.
\usepackage{minted}
\setminted{autogobble,linenos}

\usepackage{pythontex}
% The build runs pythontex from the project root (latexmk's mylatex fudge), so
% paths are relative to the project root: keep outputs here and run Python code
% from the project root.
\setpythontexoutputdir{.}
\setpythontexworkingdir{.}
\setpythontexfv{numbers=left}

\usepackage{amsmath}
\usepackage{amssymb}
\usepackage{mathtools}
\usepackage{amsthm}
\usepackage{thmtools}
\usepackage[unq]{unique}
\DeclareMathOperator{\powerset}{\mathcal{P}}

\usepackage[binary-units]{siunitx}

\usepackage[strict]{csquotes}
\usepackage[single]{acro}

\usepackage[capitalize]{cleveref}
% Beamer frames carry a `framenumber` label type; give cleveref a name for it so
% \cref to a \label'd frame works (and does not error on a cold first compile).
\crefname{framenumber}{frame}{frames}
\Crefname{framenumber}{Frame}{Frames}

% Load didactic LAST: it adapts biblatex/csquotes/cleveref (already loaded
% above) and auto-detects memoir to apply the modern Tufte layout -- wide outer
% margin, Palatino fonts, footnotes-in-margin, side-captions, \ltnote margin
% notes -- and footnote citations.
\usepackage[marginparmargin=right]{didactic}

% A semantic hypothesis environment in the style of didactic's question:
% numbered theorem in the article, block in the slides.  The translations must
% be declared before \ProvideSemanticEnv fetches them.
\DeclareTranslation{English}{Hypothesis}{Hypothesis}
\DeclareTranslation{English}{hypothesis}{hypothesis}
\DeclareTranslation{English}{hypotheses}{hypotheses}
\DeclareTranslation{English}{Hypotheses}{Hypotheses}
\ProvideSemanticEnv{hypothesis}[block]{Hypothesis}
  {hypothesis}{hypotheses}{Hypothesis}{Hypotheses}

% The research questions and hypotheses as macros, so that they can be
% restated verbatim wherever needed (they are stated in question/hypothesis
% environments in the introduction and the theory section).
\newcommand{\rqpatterns}{%
  Which patterns of variation (contrast, generalisation, fusion), if any, can
  be observed in students' recorded debugging episodes?%
}
\newcommand{\rqdeep}{%
  Is a deep approach to learning associated with generating more, and more
  complete, patterns of variation while debugging?%
}
\newcommand{\rqsuccess}{%
  Is debugging success associated with the deep approach and with the
  completeness of the generated patterns?%
}
\newcommand{\rqexperts}{%
  Which patterns of variation do expert debuggers exhibit while debugging,
  and how do the students' patterns compare to theirs?%
}
\newcommand{\hpatterns}{%
  Successful debugging episodes exhibit the patterns of variation: contrast,
  generalisation and fusion.
  Unsuccessful or prolonged episodes lack one or more of the patterns.%
}
\newcommand{\hdeep}{%
  Students with a stronger deep approach (higher R-SPQ-2F deep score)
  generate more, and more complete, patterns of variation in their debugging
  episodes.%
}
\newcommand{\hsuccess}{%
  Debugging success (fixing the bug, and doing so in fewer edit--run cycles)
  is positively associated with the deep approach and with the completeness
  of the generated patterns.%
}
\newcommand{\rqgenerative}{%
  Do students who have experienced the course material's patterns of
  variation generate more, and more complete, patterns of variation in
  their debugging than students who study on their own?%
}
\newcommand{\hgenerative}{%
  Students who use the course material generate more, and more complete,
  patterns of variation in their debugging episodes than students who do
  not---controlling for approach and prior programming experience.%
}
